\scanpage{26}
% transcription
1. $n$ catégories (strictes -- pas strictes) et $n$-foncteurs (stricts ou pas).
\inserted{Il y a des chances que les ``pas stricts'' se ramènent aux stricts}
les $n$-catégories forment une $(n+1)$-catégorie, les $\underline{\operatorname{Hom}}$ y sont des
$n$-catégories ; notion de $i$-fidélité $(0\leq i\leq n+1)$ -- la $(n+1)$-fidélité
s'appelle $n$-équivalence.

Relations entre $n$-catégories et $(n+1)$-catégories \inserted{réduites} (au moins \unsure{en|dans les} 2 sens)
Une Gr-catégorie ``est'' une 2-catégorie ayant un seul 0-objet.

2. Exemples liés aux structures semi-simpliciales.
$\left\{\begin{gathered}\text{Espaces topologiques, groupes top, anneaux top etc}\\
\text{ensembles }(\tfrac12)\text{simpliciaux }\deleted{\text{\unsure{qco}}}\text{ (de Kan ?), groupes, anneaux etc}\end{gathered}\right.$
\margin{(idem)}
\unsure{Axiomatisation|Abstractisation} : catégorie $\frac12$ simpliciale

( $A$ catégorie, $C$ \deleted{\unsure{est}} catégorie (p. ex. ens. finis tot. ordonnés non vides \ldots)
la donnée d'un foncteur $\underset{\deleted{\unsure{\rm can}}}{\underline{\operatorname{Hom}}_{\cdot}}:A^o\times A\to\widehat C$
\inserted{$(A,B)\mapsto\underline{\operatorname{Hom}}_{\cdot}(A,B)$}
équivaut à celle d'un foncteur $A^o\times A\times C^o\mathbin{\unsure{-|\to}}\mathrm{Ens}$ ou encore
$C^o\times A\xrightarrow{\underline{\operatorname{Hom}}}\widehat A\quad(\sigma,B)\mapsto\underline{\operatorname{Hom}}(\sigma,B)$ ou d'un foncteur \inserted{$\sigma\mathbin{-}\varphi_\sigma$}
$C^o\longrightarrow\underline{\operatorname{Hom}}(A,\widehat A)\bigl[\cong\underline{\operatorname{Hom}}_{!}(\widehat A,\widehat A)\bigr]$ (le $!$ indique les
foncteurs qui commutent aux $\varinjlim$ quelc.) \deleted{soit} ou
encore un foncteur $\widehat{C}^{o}\to\underline{\operatorname{Hom}}(A,\widehat A)\simeq\underline{\operatorname{Hom}}_{!}(\widehat A,\widehat A)$
commutant aux $\varprojlim$ quelc, qu'on note encore $\sigma\to h_\sigma$.

Si $C$ a objet final $e$ (exemples dans le cas \unsure{spécifié|précisé})
alors la donnée d'un isom bifonctoriel
\[
\operatorname{Hom}_{\cdot}(A,B)(e)\simeq\operatorname{Hom}(A,B)\tag*{($\alpha$)}
\]
équivaut à la donnée d'un isom.
\[
\varphi_e=\mathrm{id}_{\widehat A}\tag*{($\beta$)}
\]

La donnée d'accouplements de trifoncteurs
\[
\underset{\deleted{\unsure{\rm can}}}{\operatorname{Hom}_{\cdot}}(A,B)\times\operatorname{Hom}_{\cdot}(B,C)\mathbin{-}\operatorname{Hom}_{\cdot}(A,C)\tag*{(a)}
\]
équivaut sauf erreur à celle d'hom\inserted{om}orphismes
bifonctoriels
\[
\varphi_{K\times L}\mathrel{\overset{\deleted{\unsure{\sim}}}{\leftarrow}}\varphi_K\circ\varphi_L\tag*{(b)}
\]
\margin{Quand \inserted{(b)} est-il un isomorphisme ? sss \textcircled{\scriptsize b$'$}
$\bigl\{$\diagram{\resizebox{\linewidth}{!}{\begin{tikzcd}[ampersand replacement=\&,row sep=small]
\operatorname{Hom}_{\cdot}(A,\underline{\operatorname{Hom}}(L,B))\arrow[d]\\
\underline{\operatorname{Hom}}(L,\operatorname{Hom}_{\cdot}(A,B))
\end{tikzcd}}}est un isom.}
(ou encore, d'hom fonctoriels
\[
\varphi_M\longleftarrow\varphi_{M'}\circ\varphi_N\tag*{(c)}
\]
)
et la \deleted{donnée d'un isom} \inserted{condition} d'associativité \deleted{\unsure{est}} pour
les accouplements entre les $\operatorname{Hom}_{\cdot}$ équivaut
sauf erreur à celle d'une \deleted{\unsure{isom}} compatibilité
d'associativité pour les données (b). Compatibilité
% archive sheet number: 25
